
引理~\ref{lem:7-2-310} 的证明
将以障碍理论为基础；
下面的命题 7.2.312 概括所需结果。

设 $C'_0$、$C'$ 是滤过
$A_\infty$ 代数，
$\mathfrak h\colon C'_0\to C'$
是严格滤过 $A_{n',K'}$ 同态。
设
$\widehat{\mathfrak C}$、
$\widehat{\mathfrak C}_{s=s_0}$、
$\widehat{\mathfrak C}_{r=r_0}$，
以及
$C=(C_0,\{C_{r_0,s_0}\}_{r_0,s_0})$
如定义~\ref{def:7-2-188}。
也就是说，
$\widehat{\mathfrak C}$
是 $[0,1]^2\times C$ 的模型。
假定所有对象均为 $\mathfrak G$-有隙。

\begin{proposition}{7.2.312}\label{prop:7-2-312}
设
\[
\begin{gathered}
\mathfrak H_0\colon C'_0\longrightarrow
\widehat{\mathfrak C},\\
\mathfrak H^{r=r_0}\colon C'\longrightarrow
\widehat{\mathfrak C}_{r=r_0},\qquad
\mathfrak H^{s=s_0}\colon C'\longrightarrow
\widehat{\mathfrak C}_{s=s_0}
\end{gathered}
\]
是 $\mathfrak G$-有隙滤过
$A_{n',K'}$ 同态。
假定 $\mathfrak h$ 是同伦等价，
且 $\mathfrak h_{1,\beta_0}$ 是分裂单射，
并假定
\begin{equation}
\operatorname{Eval}_{s=s_0}\circ
\mathfrak H^{r=r_0}
=\operatorname{Eval}_{r=r_0}\circ
\mathfrak H^{s=s_0},
\tag{7.2.313.1}\label{eq:7-2-313-1}
\end{equation}
\begin{equation}
\mathfrak H^{r=r_0}\circ\mathfrak h
=\operatorname{Eval}_{r=r_0}\circ\mathfrak H_0,
\tag{7.2.313.2}\label{eq:7-2-313-2}
\end{equation}
\begin{equation}
\mathfrak H^{s=s_0}\circ\mathfrak h
=\operatorname{Eval}_{s=s_0}\circ\mathfrak H_0.
\tag{7.2.313.3}\label{eq:7-2-313-3}
\end{equation}
则存在 $\mathfrak G$-有隙滤过
$A_{n',K'}$ 同态
$\mathfrak H\colon C'\to\widehat{\mathfrak C}$，
满足
\begin{equation}
\operatorname{Eval}_{s=s_0}\circ\mathfrak H
=\mathfrak H^{s=s_0},
\tag{7.2.314.1}\label{eq:7-2-314-1}
\end{equation}
\begin{equation}
\operatorname{Eval}_{r=r_0}\circ\mathfrak H
=\mathfrak H^{r=r_0},
\tag{7.2.314.2}\label{eq:7-2-314-2}
\end{equation}
\begin{equation}
\mathfrak H\circ\mathfrak h=\mathfrak H_0.
\tag{7.2.314.3}\label{eq:7-2-314-3}
\end{equation}
\end{proposition}

须指出，
在下面的证明中，
我们完全不会使用
$\widehat{\mathfrak C}$
的 $\operatorname{Incl}$。
